%%%%%%%%%%%%%%%%%%%%%%%%%%%%%%%%%%%%%%%%%
% Masters/Doctoral Thesis 
% LaTeX Template
% Version 2.3 (25/3/16)
%
% This template has been downloaded from:
% http://www.LaTeXTemplates.com
%
% Version 2.x major modifications by:
% Vel (vel@latextemplates.com)
%
% This template is based on a template by:
% Steve Gunn (http://users.ecs.soton.ac.uk/srg/softwaretools/document/templates/)
% Sunil Patel (http://www.sunilpatel.co.uk/thesis-template/)
%
% Template license:
% CC BY-NC-SA 3.0 (http://creativecommons.org/licenses/by-nc-sa/3.0/)
%
%%%%%%%%%%%%%%%%%%%%%%%%%%%%%%%%%%%%%%%%%

%------------------------------------------------------------------------------
%	PACKAGES AND OTHER DOCUMENT CONFIGURATIONS
%------------------------------------------------------------------------------

\documentclass[
11pt, % The default document font size, options: 10pt, 11pt, 12pt
%oneside, % Two side (alternating margins) for binding by default, uncomment to switch to one side
%chapterinoneline,% Have the chapter title next to the number in one single line
%english, % ngerman for German
spanish,
singlespacing, % Single line spacing, alternatives: onehalfspacing or doublespacing
%draft, % Uncomment to enable draft mode (no pictures, no links, overfull hboxes indicated)
%nolistspacing, % If the document is onehalfspacing or doublespacing, uncomment this to set spacing in lists to single
%liststotoc, % Uncomment to add the list of figures/tables/etc to the table of contents
%toctotoc, % Uncomment to add the main table of contents to the table of contents
parskip, % Uncomment to add space between paragraphs
%nohyperref, % Uncomment to not load the hyperref package
headsepline, % Uncomment to get a line under the header
]{MastersDoctoralThesis} % The class file specifying the document structure



\usepackage[utf8]{inputenc} % Required for inputting international characters
\usepackage[T1]{fontenc} % Output font encoding for international characters

\usepackage{palatino} % Use the Palatino font by default
%,style=authoryear
\usepackage[backend=bibtex,natbib=true,sorting=none]{biblatex} % @Ale Celery
%\usepackage[backend=bibtex,natbib=true]{biblatex} % Use the bibtex backend with the authoryear citation style (which resembles APA)

\addbibresource{references.bib} % The filename of the bibliography

\usepackage[autostyle=true]{csquotes} % Required to generate language-dependent quotes in the bibliography

\usepackage{caption}
\usepackage{subcaption}

%------------------------
\usepackage{listings}

\usepackage[framemethod=tikz]{mdframed} % @Eric for License
\usepackage{enumitem} % @Eric for Lists 1. 1.1. ...

% Paquete para código en algún lenguaje de programación
\usepackage{listings} % @Eric for Syntax Highlighting

% Defino el sintaxhighlight para lenguaje C
\lstdefinestyle{customc}{
	frame=Ltb,
    framerule=0pt,
    aboveskip=0.5cm,
    framextopmargin=3pt,
    framexbottommargin=3pt,
    framexleftmargin=0.4cm,
    framesep=0pt,
    rulesep=.4pt,
    %backgroundcolor=\color{gray97},
    rulesepcolor=\color{black},
    %
    %stringstyle=\ttfamily,
    showstringspaces = false,
    basicstyle=\small\ttfamily,
    %commentstyle=\color{gray45},
    %keywordstyle=\bfseries,
    %
    numbers=left,
    numbersep=15pt,
    numberstyle=\tiny,
    numberfirstline = false,
    breaklines=true,
    %
    %     
  belowcaptionskip=1\baselineskip,
  %breaklines=true,
  %frame=L,
  xleftmargin=\parindent,
  language=C,
  %showstringspaces=false,
  %basicstyle=\footnotesize\ttfamily,
  keywordstyle=\bfseries\color{purple!80!black},
  commentstyle=\itshape\color{green!40!black},
  identifierstyle=\color{blue},
  stringstyle=\color{orange},
  %numbers=left, % where to put the line-numbers; values are (none, left, right)
  %numbersep=10pt,   % how far the line-numbers are from the code
  %numberstyle=\tiny % the style that is used for the line-numbers
}	
\lstset{escapechar=@,style=customc}	

% Defino el sintaxhighlight para lenguaje Smalltalk

\lstdefinelanguage{Smalltalk}{
  morekeywords={true,false,self,super,nil},
  sensitive=true,
  morecomment=[s]{"}{"},
  morestring=[d]',
  %style=SmalltalkStyle
}
%\lstdefinestyle{SmalltalkStyle}{
%  literate={:=}{{$\gets$}}1{^}{{$\uparrow$}}1
%} 




%\usepackage[hyphens]{url}
%\usepackage[hidelinks]{hyperref}
%\hypersetup{breaklinks=true}
\urlstyle{same}
%\usepackage{cite}

%--------------------------

\usepackage{color}

%
%------------------------------------------------------------------------------
%	MARGIN SETTINGS
%------------------------------------------------------------------------------

\geometry{
	paper=a4paper, % Change to letterpaper for US letter
	inner=2cm, % Inner margin
	outer=3.3cm, % Outer margin
	bindingoffset=2cm, % Binding offset
	top=1.5cm, % Top margin
	bottom=1.5cm, % Bottom margin
	%showframe,% show how the type block is set on the page
}

%------------------------------------------------------------------------------
%	INFORMACIÓN DE LA MEMORIA
%------------------------------------------------------------------------------

\thesistitle{sAPI (simpleAPI): diseño e implementación de una biblioteca para sistematizar
la programación de sistemas embebidos} % El títulos de la memoria, se usa en la carátula y se puede usar el cualquier lugar del documento con el comando \ttitle
\supervisor{MSc. Ing. Félix Gustavo Emilio Safar (UNQ)} % El nombre del director, se usa en la carátula y se puede usar el cualquier lugar del documento con el comando \supname
\degree{Magíster en Sistemas Embebidos} % Nombre del grado, se usa en la carátula y se puede usar el cualquier lugar del documento con el comando \degreename
\author{Esp. Ing. Eric Nicolás Pernia} % Tu nombre, se usa en la carátula y se puede usar el cualquier lugar del documento con el comando \authorname
\juradoUNO{Dr. Ing. Pablo Martín Gómez (UBA)} % Nombre y pertenencia del un jurado se usa en la carátula y se puede usar el cualquier lugar del documento con el comando \jur1name
\juradoDOS{Mg. Ing. Pablo Oscar Ridolfi (UTN FRBA)} % Nombre y pertenencia del un jurado se usa en la carátula y se puede usar el cualquier lugar del documento con el comando \jur2name
\juradoTRES{Ing. Juan Manuel Cruz (FIUBA,UTN-FRBA)} % Nombre y pertenencia del un jurado se usa en la carátula y se puede usar el cualquier lugar del documento con el comando \jur3name
\fechaINICIO{marzo de 2018}
\fechaFINAL{diciembre de 2018}

% Your subject area, this is not currently used anywhere in the template, print it elsewhere with \subjectname
\subject{Memoria del Trabajo Final de la Maestría en Sistemas Embebidos (MSE) de la UBA} 
\keywords{MSE, Sistemas Embebidos, CIAA} % Keywords for your thesis, this is not currently used anywhere in the template, print it elsewhere with \keywordnames
\university{Universidad de Buenos Aires} % Your university's name and URL, this is used in the title page and abstract, print it elsewhere with \univname
\faculty{{Facultad de Ingeniería}} % Your faculty's name and URL, this is used in the title page and abstract, print it elsewhere with \facname
\department{Departamento de Electrónica} % Your department's name and URL, this is used in the title page and abstract, print it elsewhere with \deptname
\group{{Laboratorio de Sistemas Embebidos}} % Your research group's name and URL, this is used in the title page, print it elsewhere with \groupname


\hypersetup{pdftitle=\ttitle} % Set the PDF's title to your title
\hypersetup{pdfauthor=\authorname} % Set the PDF's author to your name
\hypersetup{pdfkeywords=\keywordnames} % Set the PDF's keywords to your keywords


\newcaptionname{spanish}{\acknowledgementname}{Agradecimientos}
\newcaptionname{spanish}{\authorshipname}{Declaración de Autoría}
\newcaptionname{spanish}{\abbrevname}{Glosario}
\newcaptionname{spanish}{\byname}{por}

\renewcommand{\lstlistingname}{Algoritmo}% Listing -> Algorithm

% List of Listings -> List of Algorithms
\renewcommand{\lstlistlistingname}{Índice de \lstlistingname s}

\renewcommand{\listtablename}{Índice de Tablas}
\renewcommand{\tablename}{Tabla} 

\addtolength{\footnotesep}{2mm} % Espacio adicional en los footnotes

\begin{document}

% Use roman page numbering style (i, ii, iii, iv...) for the pre-content pages
\frontmatter 

% Default to the plain heading style until the thesis style is called for the
% body content
\pagestyle{plain}

%------------------------------------------------------------------------------
%	CARÁTULA
%------------------------------------------------------------------------------

\begin{titlepage}
\begin{center}

{\scshape\LARGE UNIVERSIDAD DE BUENOS AIRES\par}\vspace{0.1cm} % University name
{\scshape\LARGE FACULTAD DE INGENIERÍA\par}\vspace{0.1cm} % Faculty name
{\scshape\LARGE Maestría en Sistemas Embebidos\par}\vspace{1cm} % Thesis type

\includegraphics[width=.3\textwidth]{./Figures/logoFIUBA.png}
\vspace{0.7cm}

\textsc{\Large Memoria del Trabajo Final}\\[0.5cm] % Thesis type

{\huge \bfseries \ttitle\par}\vspace{0.4cm} % Thesis title

\vspace{0.7cm}
\LARGE\textbf{Autor:\\
\authorname}\\ % Author name

\vspace{1cm}

\large
\vspace{10px}
{Director:} \\
{\supname} % Supervisor name
 
\vspace{0.7cm}
Jurados:\\
\jurunoname\\
\jurdosname\\
\jurtresname
 
\vfill
\textit{Este trabajo fue realizado en las Ciudad Autónoma de Buenos Aires, entre \fechaINICIOname \hspace{1px} y \fechaFINALname.}
\end{center}
\end{titlepage}


%------------------------------------------------------------------------------
%	LICENCIA - LICENSE
%------------------------------------------------------------------------------
\newpage
%\begin{center}
%
%\begin{mdframed}[outerlinecolor=black, outerlinewidth=2pt, linecolor=gray, middlelinewidth=3pt, roundcorner=10pt]
%
% \begin{center}
%  \includegraphics*{Figures/cc-by-sa.png}
% \end{center}
%
%\emph{sAPI (simpleAPI): diseño e implementación de una biblioteca para
%sistematizar la programación de sistemas embebidos} por Esp. Ing. Eric Nicolás
%Pernia se distribuye bajo una \textbf{Licencia Creative Commons
%Atribución-CompartirIgual 4.0 Internacional}.
%Para ver una copia de esta licencia, visita:
%\\
%\url{http://creativecommons.org/licenses/by-sa/4.0/}.
%
%\end{mdframed}
%
%\end{center}


%------------------------------------------------------------------------------
%	RESUMEN - ABSTRACT 
%------------------------------------------------------------------------------

\begin{abstract}
\addchaptertocentry{\abstractname} % Add the abstract to the table of contents
%
%The Thesis Abstract is written here (and usually kept to just this page). 
%The page is kept centered vertically so can expand into the blank space above
%the title too\ldots
\centering

En esta memoria se presenta el diseño de una biblioteca para la programación de
sistemas embebidos portable entre plataformas de hardware y lenguajes de
programación. Se realizó una implementación de referencia en lenguaje C para
las plataformas del Proyecto CIAA. Se realizó un banco de pruebas de hardware,
junto a la utilización de testeo unitario e integración contínua para la
validación.

Además de la biblioteca, se creó una herramienta de código abierto para
desarrolladores que automatiza la implemetación de bibliotecas a partir de un
modelo que las describe. De esta manera se facilita la futura ampliación de la
biblioteca y su implementación en otras plataformas de hardware.

\end{abstract}

%------------------------------------------------------------------------------
%	CONTENIDO DE LA MEMORIA  - AGRADECIMIENTOS
%------------------------------------------------------------------------------

\begin{acknowledgements}
% Descomentando esta línea se puede agregar los agradecimientos al índice
%\addchaptertocentry{\acknowledgementname} 
\vspace{1.5cm}

Agradecimientos personales. \textbf{[OPCIONAL]} 

No olvidarse de agradecer al tutor.

No vale poner anti-agradecimientos (este trabajo fue posible a pesar de...)

\end{acknowledgements}

%------------------------------------------------------------------------------
%	LISTA DE CONTENIDOS/FIGURAS/TABLAS
%------------------------------------------------------------------------------
\renewcommand{\listtablename}{Índice de Tablas}

\tableofcontents % Prints the main table of contents

\listoffigures % Prints the list of figures

\listoftables % Prints the list of tables


%------------------------------------------------------------------------------
%	CONTENIDO DE LA MEMORIA  - DEDICATORIA
%------------------------------------------------------------------------------

% escribir acá si se desea una dedicatoria
\dedicatory{\textbf{Dedicado a... [OPCIONAL]}}  

%------------------------------------------------------------------------------
%	CONTENIDO DE LA MEMORIA  - CAPÍTULOS
%------------------------------------------------------------------------------

\mainmatter % Begin numeric (1,2,3...) page numbering

\pagestyle{thesis} % Return the page headers back to the "thesis" style

\renewcommand{\tablename}{Tabla} 

%------------------------------------------------------------------------------

% Define some commands to keep the formatting separated from the content 
\newcommand{\keyword}[1]{\textbf{#1}}
\newcommand{\tabhead}[1]{\textbf{#1}}
\newcommand{\code}[1]{\texttt{#1}}
\newcommand{\file}[1]{\texttt{\bfseries#1}}
\newcommand{\option}[1]{\texttt{\itshape#1}}
\newcommand{\grados}{$^{\circ}$}

%------------------------------------------------------------------------------

\newcommand{\titulo}[1]{\bigskip \noindent\textbf{#1} \\} % @Eric
\newcommand{\subtitulo}[1]{\bigskip \noindent\textit{#1} \\} % @Eric

%------------------------------------------------------------------------------

% Incluir los capítulos como archivos separados desde la carpeta Chapters
% Descomentar las líneas a medida que se escriben los capítulos

% Introducción General
\chapter{Introducción General}
\label{ChapterIntroGral}

En este capítulo se presenta el contexto y justificación de este trabajo final, la motivación para llevarlo a cabo, sus objetivos y alcance.
%(\ref{sec:contextoYJustificacion})
%(\ref{sec:motivacion})
%(\ref{sec:objetivosAlance})

%------------------------------------------------------------------------------
%	SECTION
%------------------------------------------------------------------------------
\section{Motivación}
\label{sec:motivacion}

En la actualidad existe una enorme variedad de plataformas de sistemas embebidos en el mercado, y si bien todas cuentan con dispositivos programables con características similares y periféricos compatibles, se observa que a la hora de programarlas todas son distintas. En consecuencia, el profesional desarrollador de sistemas embebidos debe invertir mucho tiempo en aprender a programar cada nueva plataforma.

Estas diferencias se deben a que los fabricantes de los dispositivos y empresas asociadas despliegan sus respectivas arquitecturas de hardware tanto en núcleos de procesamiento, como en periféricos. Cada periférico exhibe externamente la funcionalidad esperada, pero su configuración y uso es diferente. Si bien esto da a los fabricantes la posibilidad de diferenciarse frente a la competencia y a los desarrolladores elegir el dispositivo programable que más se adecue a un proyecto, también causa que el programador deba conocer en detalle la arquitectura en particular del dispositivo a la hora de utilizarlo.

Además, estas empresas en su mayoría se limitan a ofrecer información de bajo nivel para programar el hardware directamente (escribiendo registros), o bien, mediante sus propias bibliotecas escritas generalmente en lenguaje C, que están diseñadas con una gran dependencia de la arquitectura de hardware subyacente, con una limitada o nula abstracción del hardware.

Por otro lado, se están portando cada vez más lenguajes de programación a las plataformas de sistemas embebidos, que antes se reservaban para las computadoras de propósito general. Esto se debe a que las nuevas plataformas poseen cada vez más poder de procesamiento y memoria, acompañando la complejidad creciente de las aplicaciones deseadas. En consecuencia se utilizan lenguajes de programación más modernos, en parte por su mayor capacidad de abstracción, y en parte para estandarizar el lenguaje de programación de una aplicación en la que intervienen sistemas embebidos y computadoras de propósito general, siendo un ejemplo típico una aplicación de IoT (del inglés \emph{Internet of Things}), donde intervienen sistemas embebidos distribuidos, infraestructura de red y servidores.

Existen múltiples desarrollos de bibliotecas que logran una abstracción de hardware aceptable. Sin embargo, ninguna de estas se ha adoptado como estándar de facto en la industria y el alcance de las mismas se limita a cubrir plataformas de cierto ecosistema en el marco de una empresa o comunidad y soportan un único lenguaje de programación (en general C o C++).

Este trabajo contribuye a la resolución de esta problemática mediante la sistematización de la implementación de bibliotecas para sistemas embebidos. Para ello se presenta el diseño de una biblioteca para la programación de sistemas embebidos modelada independientemente del lenguaje de programación y arquitectura del hardware. 

Asimismo, se realizaron herramientas de código abierto para desarrolladores, que automatizan gran parte de la implementación de la biblioteca en una nueva plataforma en base a una descripción de la misma, así como la definición de nuevos módulos de bibliotecas. De esta manera se facilita la futura ampliación de la biblioteca y su extensión a diferentes plataformas de hardware.		

Utilizando estas herramientas se realizó una implementación de referencia en lenguaje C para las plataformas del Proyecto CIAA \citep{CIAA}.
%\footnote{Siglas de Computadora Industrial Abierta Argentina.} []. 
%Para su validación, se realizó un banco de pruebas de hardware, junto a la utilización de testeo unitario e integración continua.
Para su validación, se desarrollaron ejemplos funcionales y se utilizaron \emph{tests} unitarios e integración continua.

%------------------------------------------------------------------------------
%	SECTION
%------------------------------------------------------------------------------
\section{Motivación}
\label{sec:motivacion}


%------------------------------------------------------------------------------
\subsection{Lenguajes de programación y Hardware en Sistemas Embebidos}

Saraza...


%------------------------------------------------------------------------------
\subsection{Bibliotecas para microcontroladores ofrecidas por los fabricantes}

Saraza...

%------------------------------------------------------------------------------
\subsection{Proyecto CIAA}

De esta forma, es posible programarlos sin necesidad de conocer detalles sobre
la arquitectura subyacente. Este diseño promueve y permite la programación 
independiente del hardware, reduciendo la complejidad general del desarrollo de
sistemas embebidos. Esta biblioteca se utilizará para programar las diferentes
plataformas que componen el Proyecto Computadora Industrial Abierta Argentina 
(CIAA).


Saraza...



En el mercado se encuentra una gran variedad de plataformas basadas en microcontrolador y aunque todas ellas poseen microcontroladores con características similares y periféricos compatibles, sin embargo, sus bibliotecas son muy diferentes. Esto se debe a que cada fabricante y/o empresas asociadas ofrece sus propias bibliotecas en lenguaje C las cuales están diseñadas fuertemente dependientes de cada arquitectura de cada microcontrolador que estas plataformas contienen.
Existen también muchas bibliotecas que logran una buena abstracción del hardware en varias plataformas, pero ninguna se ha adoptado como estándar general. Esto se debe a múltiples causas, entre ellas:


En el sector de la industria automotriz existe un estándar llamado AUTOSAR[] para la estandariza- ción de la  arquitectura de sistemas electrónicos que aún no ha logrado extenderse a otras industrias y cuya definición de bibliotecas propuesta es muy extensa y compleja de implementar y también para aprender a utilizar.
Bibliotecas de drivers basadas en POSIX[] (que podemos hallar en sistemas con en Linux Embebido como Raspberry Pi[]). Su abstracción ha sido muy útil en la estandarización de drivers para PCs con sistema operativo Unix-compatible. Sin embargo, es tan alejada del hardware físico que en la práctica provoca una muy baja utilización por parte de los profesionales electrónicos y afines que se dedican a la programación de sistemas embebidos.

Programadores hobbistas de sistemas embebidos han impulsado la  estandarización de la programación mediante una biblioteca conocida como Wiring[] disponible para múltiples plataformas (como la popular Arduino[]). Esta biblioteca si bien logra una gran facilidad de uso y rápido aprendizaje contiene algunas imprecisiones técnicas que provoca vicios indeseados en el aprendizaje de programación de microcontroladores. También carece de definición de una API para la utilización del periférico temporizador el cual es muy utilizado en un microcontrolador.

En la actualidad existen muchos fabricantes de microcontroladores que adquieren licencias para la fabricación de microcontroladores con núcleos de procesamiento de arquitecturas Cortex M[] diseñados por la empresa ARM[]. Para los mismos existe una biblioteca estándar llamada CMSIS{} para la programación del núcleo de procesamiento y controlador de interrupciones pero que no se extiende a los perifŕicos donde cada fabricante busca diferenciarse de sus competidores.

Otras empresas muy difundidas en el campo de la enseñanza en la programación como Microchip[] proveen bibliotecas muy dependientes del hardware y generadores automatizados de código que si bien en principio aceleran los tiempos de desarrollo, cuando una aplicación requiere modos más avanzados terminan dificultando la programación pues unas configuraciones pisan a las otras provocando que no funcione.

De los ejemplos anteriores se observa que para la realización de una biblioteca estándar que satisfaga a los diferentes usuarios se debe lograr un balance entre:

\begin{itemize}
\item Extensión de la definición de la API.
\item Dependencia del hardware.
\item Nivel de abstracción.
\item Complejidad de aprendizaje y uso.
\item Periféricos y modos soportados.
\item Escalabilidad.
\end{itemize}



Por estos motivos se decide realizar la definición de una API de una biblioteca estándar para microcontroladores  en lenguaje C que supere todas estas dificultades.

%------------------------------------------------------------------------------
%	SECTION
%------------------------------------------------------------------------------
\section{Objetivos y alcance}
\label{sec:objetivosAlance}

En esta sección se definen los objetivos del presente Trabajo Final (sección \ref{subsec:objetivos}) y el alcance (sección \ref{subsec:alcance}) 

De esta forma, es posible programarlos sin necesidad de conocer detalles sobre
la arquitectura subyacente. Este diseño promueve y permite la programación 
independiente del hardware, reduciendo la complejidad general del desarrollo de
sistemas embebidos. Esta biblioteca se utilizará para programar las diferentes
plataformas que componen el Proyecto Computadora Industrial Abierta Argentina 
(CIAA).

\subsection{Objetivos}
\label{subsec:objetivos}

El objetivo de este proyecto es diseñar e implementar una biblioteca de
software para la programación de sistemas embebidos basados en
microcontroladores con las siguientes características: 

\begin{itemize}
   \item Estar modelada independientemente de los lenguajes de programación.
   \item Definir una interfaz de programación de aplicaciones (API) sencilla que abstraiga los modos de uso más comunes de los periféricos típicos que hallados en cualquier microcontrolador del mercado. 
   \item Ser totalmente portable entre diferentes arquitecturas de hardware sobre donde se ejecuta, manteniendo una API uniforme a lo largo de las mismas\footnote{Debe cumplir la función de capa de abstracción de hardware, o \textit{Hardware Abstraction Layer} (HAL),  en inglés}.
\end{itemize}

Dicha biblioteca se deberá implementar en lenguaje C para las plataformas del
Proyecto CIAA.

\subsection{Alcance}
\label{subsec:alcance}

Este Trabajo Final incluye realización de:

\begin{itemize}
   \item Diseño de la biblioteca. Archivos de descripción de la biblioteca mediante diferentes diagramas y código independiente del lenguaje de programación.
   \item Implementación en lenguaje C de la biblioteca para las plataformas de hardware:
   \begin{itemize}
      \item CIAA-NXP.
      \item EDU-CIAA-NXP.
      \item CIAA-Z3R0.
      \item PicoCIAA.
   \end{itemize}
   \item Manual de instalación de las herramientas para utilizar la biblioteca con las plataformas de hardware citadas.
   \item Manual de referencia de la biblioteca.
   \item Ejemplos de utilización.
\end{itemize}



% Introducción Específica
\chapter{Introducción Específica}
\label{ChapterIntroEsp}

En este capítulo se explican los detalles para comprender las decisiones de diseño adoptadas. Se describen los trabajos previos que se tomaron como fundamento para el diseño de la biblioteca, se presentan las plataformas sobre las cuales se elige realizar la implementación de la misma y los requerimientos de este trabajo, junto a la planificación para su cumplimiento.

%En particular, en la sección \ref{sec:antecedentes}
%en la sección \ref{sec:ciaaBoards}
% así como las herramientas provistas por los fabricantes de sus respectivos microcontroladores. Considerando las secciones anteriores, 
%cumplimiento de los mismos (\ref{sec:planificacion}).

%------------------------------------------------------------------------------
%	SECTION
%------------------------------------------------------------------------------
\section{Antecedentes}
\label{sec:antecedentes}

Los primeros pasos para la idea de la realización de este trabajo surgen del trabajo final de la Carrera de Especialización en Sistemas Embebidos de la Universidad de Buenos Aires (CESE FIUBA) \cite{CESE}, realizado por el autor, titulado "Desarrollo de Firmware y Software para programar la CIAA en lenguaje JAVA con aplicación en entornos Industriales" \cite{CeseTesisEric} presentado en diciembre de 2015, donde como parte del mismo se implementó la programación de plataformas de sistemas embebidos en lenguaje Java \cite{Java}. 

En ese trabajo se llevó a cabo, entre otras cosas, la implementación de una biblioteca para acceder a los periféricos de un microcontrolador en lenguaje Java. Para llevarlo a cabo se realizó la definición de clases que modelan los periféricos GPIO, ADC y UART y se implementó el acceso al hardware como métodos nativos de Java, escritos en lenguaje C. Por este motivo se debió realizar una implementación de la biblioteca también lenguaje C. 
%Esta implementación logra un buen balance entre facilidad de uso, extensión y modos de uso definidos. 
Se implementó para las plataformas del proyecto CIAA EDU-CIAA-NXP \cite{EDUCIAA} y CIAA-NXP \cite{CIAANXP}.
% (plataformas que se describen en detalle en la sección \ref{sec:ciaaBoards}). 
Esta biblioteca se denominó sAPI, siglas de \emph{simple API} en alusión a que provee una API sencilla para la programación de microcontroladores.

%Desde su desarrollo, se ha utilizado como ejemplo de capa de abstracción de hardware en las asignaturas dictadas por el autor ''Programación de microprocesadores'' [] de la CESE FIUBA, y la asignatura ''Sistemas Digitales'' de la carrera Ingeniería en Automatización y Control Industrial (IACI[]) de la Universidad Nacional de Quilmes (UNQ).

A partir mazo de 2016 se decide utilizar la biblioteca como base para la enseñanza de la programación de periféricos de microcontroladores.
% en los Cursos Abiertos de Programación de Sistemas Embebidos (CAPSE[]) organizados por la ACSE. 
A lo largo de ese año se extendió la biblioteca de C para la plataforma EDU-CIAA-NXP de forma considerable para explicar la utilización los periféricos típicos de microcontroladores, logrando excelentes resultados en aprendizaje por parte de los alumnos tanto de niveles avanzados como quienes dan sus primeros pasos en el aprendizaje de programación de microcontroladores.

Además, la biblioteca sAPI en lenguaje C para la EDU-CIAA-NXP se puso a disposición de cualquier persona ya que se encuentra publicada de forma libre y gratuita por internet bajo una licencia BSD modificada \citep{BSD3clause} en el sitio de \emph{GitHub} del autor \citep{sAPIgit}.

Finalmente, en diciembre de 2016 se decide utilizar la biblioteca sAPI realizada en lenguaje C como biblioteca estándar para las plataformas del proyecto CIAA, distribuyéndola como parte del \emph{framework}\footnote{En el desarrollo de software, un \emph{framework} es una estructura conceptual y tecnológica de asistencia definida, normalmente, con artefactos o módulos concretos de software, que sirve de base para la organización y desarrollo de software. En este caso concreto se compone de bibliotecas de código C y makefiles para su compilación permitiendo organizar proyectos de software en lenguaje C.} ''Firmware v2'' \citep{ciaaFirmwareV2}. Este \emph{framework} combinó la biblioteca sAPI con \emph{framework} "Workspace" \citep{ws-ridolfi}, desarrollado por Pablo Ridolfi.

Esta plataforma ha sido adoptada por una gran cantidad de usuarios. Una prueba de ello es la encuesta realizada en octubre de 2018, titulada: "Tecnologías usadas en los Trabajos Finales del Posgrado en Sistemas Embebidos: 2015-2018" \citep{EncuestaCeseMse} donde en la sección 27, \emph{''Uso de material del Proyecto CIAA en los trabajos finales"} se observa que alrededor de la mitad de los trabajos finales de la CESE/MSE, utilizaron material generado en el marco del Proyecto CIAA. Y en particular, a partir de 2017 se produjo un cambio de tendencia y más del 60\% de los trabajos finales utilizaron material del Proyecto CIAA. Este cambio en la tendencia se observa que comienza en 2016 con la publicación de firmware v2 como se muestra en la gráfica de la figura \ref{fig:tendenciaMaterialCiaa}, extraída de dicho artículo.

\begin{figure}[!htbp]
\begin{center}  % [width=14cm,height=8cm] [width=\textwidth]
\includegraphics*[width=12cm]{Figures/TendenciaMaterialCiaa.png}
\par\caption{Evolución en el uso de material del Proyecto CIAA.}\label{fig:tendenciaMaterialCiaa}
\end{center}
\end{figure}

A principios de 2017, el autor llevó a cabo una profunda revisión y mejora de la biblioteca de C con el objetivo de extenderla a las demás plataformas del proyecto CIAA.
% En este rediseño colaboró Martín Ribelotta. 
Ese trabajo fue compilado en un artículo y publicado en el Congreso Argentino de Sistemas Embebidos (CASE) \citep{paperSapiCASE2017} en agosto de 2017. Siendo un antecedente fundamental para la realización del presente trabajo final, porque se realizó un estudio exhaustivo de las bibliotecas existentes en el mercado; dando como resultado los siguientes puntos a considerar en el diseño de una biblioteca de C para la programación de sistemas embebidos:

\begin{itemize}
\item
Extensión de la definición de la API.
\item
Dependencia del hardware.
\item
Nivel de abstracción.
\item
Complejidad de aprendizaje y uso.
\item
Periféricos y modos soportados.
\item
Escalabilidad.
\end{itemize}

Parte de este rediseño se aplicó a la biblioteca sAPI en lenguaje C para la EDU-CIAA-NXP.

En 2017 el autor desarrolló una plataforma de hardware más económica que la EDU-CIAA-NXP, nombrada ''CIAA-Z3R0'' \cite{CIAAZ3R0}, que salió al mercado en noviembre de ese año. En diciembre el autor publicó la primera versión de la biblioteca sAPI en lenguaje C para esta plataforma dándole soporte a algunos periféricos de la misma.
% (se descibe en detalle en la sección \ref{sec:CIAA-Z3R0})

Tomando la experiencia adquirida a lo largo de estos años, el autor presenta en esta memoria un diseño mejorado para la biblioteca, ampliando el mismo e independizándolo del lenguaje de programación. Asimismo, teniendo en cuenta las tareas repetitivas que requieren la implementación de una biblioteca para diferentes plataformas de sistemas embebidos se decide desarrollar herramientas para automatizar el proceso donde resulta posible.

 
%------------------------------------------------------------------------------
%	SECTION
%------------------------------------------------------------------------------
\section{Plataformas del Proyecto CIAA}
\label{sec:ciaaBoards}



En sus comienzos el proyecto CIAA desarrolla una computadora industrial basada en un microcontrolador NXP LPC4337 (JDB 144) Dual-core Cortex-M4+Cortex-M0 a 204MHz, con 1 MB de memoria Flash y 136 KB de memoria SRAM nombrada CIAA-NXP [ ]. En base a este diseño los integrantes del proyecto CIAA diseñan una versión educativa sin las protecciones industriales, nombrada EDU-CIAA-NXP [ ]. En la figura [ ] se presenta esta plataforma.

Esta plataforma incluye todos los periféricos típicos que podemos encontrar en los microcontroladores disponibles en el mercado permitiendo la enseñanza con herramientas modernas.

Mediante la colaboración de la Red Universitaria de Sistemas Embebidos (RUSE) [ ] se han distribuido en 2015 entre 10 y 40 placas en Universidades de Argentina con carreras afines a la electrónica.

Utilizando la EDU-CIAA-NXP como plataforma base para la enseñanza se ha desarrollado, o colaborado en el desarrollo de diferentes herramientas de software y hardware abiertas, las cuales se describen en las siguientes secciones.

--------------------------------------


El proyecto de la Computadora Industrial Abierta Argentina (CIAA) nació en 2013 como una iniciativa conjunta entre el sector académico y el industrial, representados por la ACSE\footnote{Asociación Civil para la investigación, promoción y desarrollo de los Sistemas electrónicos Embebidos. Sitio web: \url{http://www.sase.com.ar/asociacion-civil-sistemas-embebidos}} y CADIEEL\footnote{Cámara Argentina de Industrias Electrónicas, Electromecánicas y Luminotécnicas. Sitio web: \url{http://www.cadieel.org.ar/}}, respectivamente.

\medskip

\noindent Los objetivos del proyecto CIAA son:

\begin{itemize}
\item
Impulsar el desarrollo tecnológico nacional, a partir de sumar valor agregado al trabajo y a los productos y servicios, mediante el uso de sistemas electrónicos, en el marco de la vinculación de las instituciones educativas y el sistema científico-tecnológico con la industria.
\item
Darle visibilidad positiva a la electrónica argentina.
\item
Generar cambios estructurales en la forma en la que se desarrollan y utilizan en nuestro país los conocimientos en el ámbito de la electrónica y de las instituciones y empresas que hacen uso de ella.
\end{itemize}

\noindent Todo esto en el marco de un trabajo libre, colaborativo y articulado entre industria y academia. 

\bigskip

Con esta iniciativa, se han desarrollado en la actualidad varias plataformas de hardware y entornos de programación para utilizarlas.

\medskip

Al momento de la presentación de este trabajo, existen dos versiones de la plataforma CIAA cuyo desarrollo ha sido completado:

\begin{itemize}
\item 
CIAA-NXP, basada en el microcontrolador NXP LPC4337, que ya se comercializa.
\item 
CIAA-FSL, que utiliza, en cambio, el microcontrolador Freescale MK60FX512VLQ15, pero únicamente hay prototipos de esta plataforma. 
\end{itemize}

Además, existe una versión educativa de bajo costo de la CIAA-NXP, nombrada EDU-CIAA-NXP, que ya se distribuyeron alrededor de 1000 unidades y ya hay otras 1000 reservadas en producción. 

\medskip

Debido a estas razones, el trabajo se enfoca en el desarrollo de herramientas para programar las dos plataformas basadas en el microcontrolador NXP LPC4337. Se introducen a continuación las características de las mismas.




Siendo el autor participante de este proyecto desde mediados de 2014, ocupando el rol de Responsable de Software-PLC mediante el aporte al proyecto CIAA de un IDE\footnote{IDE4PLC. Sitio web: \url{http://proyecto-ciaa.com.ar/devwiki/doku.php?id=desarrollo:software-plc}} que permite programar esta plataforma con lenguajes de PLC industriales (IEC-661131-3), se desea agregar en esta oportunidad la posibilidad de programar a esta plataforma con un lenguaje de programación orientado a objetos mediante el desarrollo de un IDE para tal fin.


\subsection{CIAA-NXP}

La CIAA-NXP es la primera y única computadora del mundo que reúne dos cualidades:

\begin{itemize}
\item 
Ser \textbf{Industrial}, ya que su diseño está preparado para las exigencias de confiabilidad, temperatura, vibraciones, ruido electromagnético, tensiones, cortocircuitos, etc., que demandan los productos y procesos industriales.
\item 
Ser \textbf{Abierta}, ya que toda la información sobre su diseño de hardware, firmware, software, etc. está libremente disponible en Internet bajo la Licencia BSD, para que cualquiera la utilice como quiera.
\end{itemize}

\noindent Esta plataforma se compone de:

\begin{itemize}
\item
CPU: Microcontrolador NXP LPC 4337 JDB 144 (Dual-core Cortex-M4 + Cortex-M0 @ 204MHz).
\item
Debugger: USB-to-JTAG FT2232H. Soportado por OpenOCD.
\item
Memorias: 
   \begin{itemize}
   \item
   IS42S16400F - SDRAM. 64Mbit @ 143MHz.
   \item
   S25FL032P0XMFI011 - Flash SPI. 32 Mbit, Quad I/O Fast read: 80 MHz.
   \item
   24AA1025 - EEPROM I2C. 1 Mbit, 400 kHz. Almacenamiento de propósito general, datos de calibración del usuario, etc.
   \item
   24AA025E48 - EEPROM I2C. 2 kbit, 400 kHz. Para implementación de MAC-Address o almacenamiento de propósito general.
   \end{itemize}
\item
Entradas y salidas:   
   \begin{itemize}
   \item
   8 entradas digitales opto-aisladas 24VDC.
   \item
   4 Entradas analógicas 0-10V/4-20mA.
   \item
   4 salidas Open-Drain 24VDC.
   \item
   4 Salidas con Relay DPDT.
   \item
   1 Salida analógica 0-10V/4-20mA.
   \end{itemize}
\item
LV-GPIO:
   \begin{itemize}
   \item
   14 GPIOs.
   \item
   I2C.
   \item
   SPI.
   \item
   4 canales analógicos.
   \item
   Aux. USB.
   \end{itemize}
\item
Interfaces de comunicación:
   \begin{itemize}
   \item
   Ethernet.
   \item
   USB On-The-Go.
   \item
   RS232.
   \item
   RS485.
   \item
   CAN.
   \end{itemize}
\item
Múltiples fuentes de alimentación.

\end{itemize}

\noindent En la figura [\ref{fig:ciaaNxp}] se muestra una fotografía de la plataforma.

\begin{figure}[!htbp]
\begin{center}  %[width=14cm,height=8cm]
\includegraphics*[width=\textwidth]{Figures/CIAA-NXP_Foto.png}
\par\caption{Plataforma CIAA-NXP.}\label{fig:ciaaNxp}
\end{center}
\end{figure}


%http://proyecto-ciaa.com.ar/devwiki/doku.php?id=desarrollo:hardware:ciaa_nxp:ciaa_nxp_inicio

\subsection{EDU-CIAA-NXP}

La plataforma EDU-CIAA-NXP es un desarrollo colaborativo, realizado por miembros de la Red Universitaria de Sistemas Embebidos (RUSE), en el marco del Proyecto CIAA. RUSE se compone de docentes pertenecientes a más de 60 Universidades a lo largo y a lo ancho del país. 

\medskip

\noindent Los propósitos de la plataforma son: 

\begin{itemize}
\item
Proveer una plataforma de desarrollo moderna, económica y de fabricación nacional basada en la CIAA-NXP, que sirva a docentes y a estudiantes en los cursos de sistemas embebidos.
\item
Lograr una amplia inserción en el sistema educativo argentino.
\item
Realizar un aporte eficaz al desarrollo de vocaciones tempranas en electrónica, computación e informática.
\item
Demostrar que las universidades argentinas son capaces de realizar un desarrollo colaborativo exitoso en el área de los sistemas embebidos, cumpliendo con requerimientos de tiempo y forma.
\end{itemize}

\medskip

\noindent Características de la EDU-CIAA-NXP:

\begin{itemize}
\item
CPU: Microcontrolador NXP LPC 4337 JDB 144 (Dual-core Cortex-M4 + Cortex-M0 @ 204MHz).
\item
Debugger: USB-to-JTAG FT2232H. Soportado por OpenOCD.
\item
2 puertos micro-USB (uno para aplicaciones y debug, otro OTG).
\item
6 salidas digitales implementadas con leds (3 normales y uno RGB).
\item
4 entradas digitales con pulsadores.
\item
1 puerto de comunicaciones RS-485 con bornera.
\item
2 conectores de expasión:
   \begin{itemize}
   \item
   P0:
      \begin{itemize}
      \item
      3 entradas analógicas (ADC0 a ADC2).
      \item
      1 salida analógica (DAC0).
      \item
      1 conexión para un teclado de 3 x 4.
      \item
      12 pines genéricos de I/O.
      \end{itemize}
   \item
   P1:
      \begin{itemize}
      \item
      1 puerto Ethernet.
      \item
      1 puerto CAN.
      \item
      1 puerto SPI.
      \item
      1 puerto I2C.
      \item
      12 pines genéricos de I/O.
      \end{itemize}
   \end{itemize}
\end{itemize}

\noindent En la figura [\ref{fig:eduCiaa}] se muestra una fotografía de esta plataforma.

\begin{figure}[!htbp]
\begin{center}  %[width=14cm,height=8cm]
\includegraphics*[width=\textwidth]{Figures/EDU-CIAA-NXP_Foto.png}
\par\caption{Plataforma EDU-CIAA-NXP.}\label{fig:eduCiaa}
\end{center}
\end{figure}

\medskip 

%http://proyecto-ciaa.com.ar/devwiki/doku.php?id=desarrollo:edu-ciaa:edu-ciaa-nxp

% http://proyecto-ciaa.com.ar/devwiki/doku.php?id=start

\subsection{Pico-CIAA}

http://www.proyecto-ciaa.com.ar/devwiki/doku.php?id=desarrollo:hardware:picociaa

Procesador Dual-Core de 32 bits

Tiene un procesador Dual Core de 32-bit ARM Cortex-M4F / Cortex-M0+ @100 MHz con 512kB flash y 104 kB SRAM, incluye debug vía un LPC11U35 programable.

Diagrama en bloques

Posee diversos puertos de comunicación, entre ellos USB, PCIExpress, UART, SPI, I2C y soporte para PWM y entradas y salidas digitales de propósito general.
special

Es la CIAA más pequeña

Con sólo 51 x 30 mm la picoCIAA es la placa más pequeña de la familia CIAA,
lo que la hace ideal para aplicaciones de internet de las cosas (IoT).

Ideal en Single Board Computers

Su interfaz PCI Express permite integrarla fácilmente en SBC de propósito general, para dar soporte de bajo nivel al manejo de dispositivos en tiempo real.

\begin{figure}[!htbp]
\begin{center}  %[width=14cm,height=8cm]
\includegraphics*[width=\textwidth]{Figures/PicoCIAA_Foto.png}
\par\caption{Plataforma PicoCIAA.}\label{fig:picoCiaa}
\end{center}
\end{figure}

\subsection{CIAA-Z3R0}

Esta plataforma es ideal para aplicaciones de bajo consumo y proyectos de robótica educativa. Está diseñada para ser integrada como componente en un diseño mayor debido a su reducido tamaño conectándola mediante tiras de pines, o soldada a través de su borde de agujeros para montaje castellated.

Se puede comprar en Argentina por aproximadamente AR \$400 a inicios de 2018.

Posee la mayoría de los periféricos que se encuentran en la EDU-CIAA-NXP pero utiliza un microcontrolador siete veces más económico que esta última, de la empresa Silicon Labs, modelo EFM32HG322F64 (QFP48) con núcleo ARM Cortex-M0+ a 25MHz, 64KB de memoria Flash y 8KB de memoria SRAM, que es suficiente para que un alumno entre en el mundo de los microcontroladores modernos de 32 bits.

El dispositivo de depuración se debe comprar por separado, sin embargo, mediante un único dispositivo de depuración se pueden programar muchas plataformas CIAAZ3R0.
No es necesario dejar este circuito de depuración en el diseño final

\begin{figure}[!htbp]
\begin{center}  %[width=14cm,height=8cm]
\includegraphics*[width=\textwidth]{Figures/CIAA-Z3R0_Foto.png}
\par\caption{Plataforma CIAA-Z3R0.}\label{fig:ciaaZero}
\end{center}
\end{figure}

\subsection{Arquitectura de Hardware}

Saraza...

\subsection{Bibliotecas disponibles}

Saraza...
 
%------------------------------------------------------------------------------
%	SECTION
%------------------------------------------------------------------------------
\section{Requerimientos}
\label{sec:requerimientos}

%\subsection{Diseño de la biblioteca}

%\subsection{Implementación}

%\subsection{Documentación y difusión}

Los requerimientos se establecieron en base a los objetivos expuestos en la sección \ref{sec:objetivosAlance}, reuniones con desarrolladores del Proyecto CIAA, alumnos (de grado de UNQ, posgrado de FIUBA y cursos CAPSE) y el director. El espíritu de los mismos es proponer una solución que permita la estandarización de la programación de microcontroladores. Estos son:

\begin{enumerate}
   \item Fecha de finalización: 19/11/2018.
   \item Diseño de la biblioteca.
      \begin{enumerate}[1]
         \item Realizar un diseño independiente del hardware y lenguaje de programación, teniendo en cuenta los conceptos familiares al programador de Sistemas Embebidos.
         \item Debe estar modelada con objetos y contar con una descripción mediante diagramas UML.
         \item Debe modelar al menos:
            \begin{enumerate}[1]
               \item CORE: un núcleo de procesamiento.
               \item GPIO: periférico que consiste en un único pin de entrada/salida de propósito general (pin), así como un grupo de pines (port).
               \item ADC: periférico conversor analógico-digital.
               \item DAC: periférico conversor digital-analógico.
               \item TIMER: periférico temporizador.CORE: un núcleo de procesamiento.
               \item RTC: periférico reloj de tiempo real.CORE: un núcleo de procesamiento.
               \item UART: periférico de comunicación serial asincrónico.
               \item SPI: periférico interfaz serie sincrónica.
               \item I2C: periférico de comunicación serie entre circuitos integrados.
            \end{enumerate}
      \end{enumerate}
   \item Implementación de la biblioteca.
      \begin{enumerate}[1]
         \item Utilizar un sistema de control de versiones con repositorios on line.
         \item Programar en lenguaje C la biblioteca para cada plataforma de hardware particular utilizando como plantilla los archivos generados.
         \item Desarrollar ejemplos de utilización para las diferentes plataformas.
      \end{enumerate}
   \item Documentación y difusión.
      \begin{enumerate}[1]
         \item Confeccionar un manual de referencia de la API de la biblioteca.
         \item Desarrollar un tutorial de instalación de las herramientas para utilizar la biblioteca con las plataformas de hardware citadas.
         \item Publicación on line del código fuente.
         \item Informar a la comunidad del Proyecto CIAA y a la comunidad de programadores de
   Sistemas Embebidos.
      \end{enumerate}
\end{enumerate}

% \subsection{Entregables principales del trabajo}
% 
% Se listan a continuación los entregables principales del proyecto:
% 
% 1. Archivos de descripción de la biblioteca mediante diferentes diagramas y código
% independiente del lenguaje de programación.
% 2. Código fuente de la implementación en lenguaje C de la biblioteca para las plataformas de hardware:
%    - CIAA-NXP.
%    - EDU-CIAA-NXP.
%    - CIAA-Z3R0.
%    - PicoCIAA.
% 3. Ejemplos de utilización.
% 4. Informe final.
% 5. Manual de referencia de la API de la biblioteca.
% 6. Tutorial de instalación de las herramientas para utilizar la biblioteca con las plataformas de hardware citadas.

 
%------------------------------------------------------------------------------
%	SECTION
%------------------------------------------------------------------------------
\section{Planificación}
\label{sec:planificacion}

En la tabla \ref{tab:planificacion} se resume la planificación del trabajo. En la misma se pueden observar cada una de las tareas planeadas, junto a su duración, fecha de inicio y fecha de finalización estimadas.

\begin{table}[!htbp]
\centering
\includegraphics*[width=16cm]{Figures/GanttSP.png}
\par\caption{Planificación del Trabajo Final.}\label{tab:planificacion}
\end{table} 

% Diseño
\chapter{Diseño e implementación}
\label{ChapterDiseno}

En este capítulo se presenta el diseño de la biblioteca sAPI con un enfoque \emph{top-down}. Se expone la descripción general del desarrollo. Se presenta el diseño de un modelo de plataforma de hardware y módulos de biblioteca sAPI, junto a su verificación. Se describe la implementación del código C dependiente del hardware y se presenta el diseño e implementación del generador de código de la biblioteca.

%------------------------------------------------------------------------------
%	SECTION
%------------------------------------------------------------------------------
\section{Descripción general}
\label{sec:descripGralDiseno}

El sistema propuesto para la creación de bibliotecas de software para la programación de plataformas de hardware se ilustra en la figura \ref{fig:sapi_gen}.

\begin{figure}[!htbp]
\begin{center}  % [width=14cm,height=8cm] [width=\textwidth]
\includegraphics*[width=14cm]{Figures/sapi_gen.pdf}
\par\caption{Diagrama del sistema diseñado.}\label{fig:sapi_gen}
\end{center}
\end{figure}

Se diseñó un modelo de plataforma de hardware (en adelante \emph{board} por sus siglas en inglés). Este modelo describe una plataforma de hardware completa y se debe instanciar para cada nueva plataforma de hardware.

Por otra parte, se desarrollaron módulos de biblioteca (en adelante módulos sAPI) independientes del hardware y lenguajes de programación. Estos definen propiedades y métodos describen las propiedades y métodos de cada uno de los periféricos de un cierto SoC, definiendo entonces, la API de la biblioteca.

Además, se debe proveer la implementación de cada uno de los métodos de los módulos sAPI (en adelante \emph{drivers}). Estos \emph{drivers} deben estar escritos en un cierto lenguaje de programación y son dependientes de la arquitectura del hardware.

También, se diseñó un generador de biblioteca sAPI. Este generador se debe definir para cada lenguaje de programación. Para este trabajo final se desarrolló solamente el generador de lenguaje C, sin embargo, el mismo se realizó de forma modular para que pueda ser fácilmente adaptado a otros lenguajes. 

Mediante el generador de C, utilizando una instancia de \emph{board} y \emph{drivers} específicos, junto con los módulos sAPI, se genera entonces una biblioteca sAPI en lenguaje C para una plataforma particular.

En las siguientes secciones de detalla el diseño e implementación de cada una de las entidades descriptas. Las mismas se modelaron utilizaron los conceptos del paradigma de la programación orientada a objetos, que son especialmente útiles para describir módulos de software que encapsulan funcionalidad. Estos se describen mediante diagramas UML de clases de forma independiente del lenguaje de programación.


% Board
%------------------------------------------------------------------------------
\subsection{Plataforma y componentes}



% Modulos sAPI
%------------------------------------------------------------------------------
%	SECTION
%------------------------------------------------------------------------------
\section{Modelo abstracto de módulo de biblioteca}
\label{sec:modelLibrary}

Cada módulo que compone la biblioteca sAPI se modela utilizando los conceptos del paradigma de la programación orientada a objetos, los cuales son especialmente útiles para describir módulos de software que encapsulan funcionalidad. Estos se describen mediante diagramas UML de clases de forma independiente del lenguaje de programación.

 

% Verificacion del modelo
%------------------------------------------------------------------------------
%	SECTION
%------------------------------------------------------------------------------
\section{Verificación del modelo}
\label{sec:modelVerification}

%------------------------------------------------------------------------------
%\subsection{}



% Codigo C dependiente de la plataforma
%------------------------------------------------------------------------------
%	SECTION
%------------------------------------------------------------------------------
\section{Implementación del código C dependiende del hardware}
\label{sec:codeImplementation}

%------------------------------------------------------------------------------
%\subsection{}

En la implementación para las plataformas del proyecto CIAA, del código dependiente del hardware, que forma parte de la biblioteca sAPI, se utilizaron las bibliotecas de drivers provistas por los fabricantes. En particular, para los microcontroladores de la empresa NXP se utilizó LPCOpen, en la version 3.02 para el LPC4337 y versión 3.04 para el LPC54102. Para el microcontrolador EFM32HG322 de Siicon Labs se utilizó la biblioteca EMLIB versión 5.1.2 (todas en sus últimas versiones al momento de la realización de este trabajo).

Por otra parte, se reutilizó parte del código desarrollado en versiones anteriores de la biblioteca sAPI para las plataformas EDU-CIAA-NXP y CIAA-Z3R0 y parte del código de PicoFirmware \citep{PicoAPI} para la implementación en la plataforma PicoCIAA.

Una aplicación de embebidos típica que utiliza la biblioteca sAPI, se puede combinar con un sistema operativo de tiempo real, \emph{stracks} y \emph{middelware} resultando una arquitectura de capas de software típica en aplicaciones de sistemas embebidos como se ilustra en la figura \ref{fig:sapiCapas2}.

\begin{figure}[!htbp]
\begin{center}  % [width=14cm,height=8cm] [width=\textwidth]
\includegraphics*[width=10.4cm]{Figures/sapiCapas2.png}
\par\caption{Arquitectura de una aplicación que utiliza la biblioteca.}\label{fig:sapiCapas2}
\end{center}
\end{figure}


% Generador de sAPI
%------------------------------------------------------------------------------
%	SECTION
%------------------------------------------------------------------------------
\section{Generador de sAPI}
\label{sec:modelGenerator}

El generador de sAPI toma como entrada el una instancia de \emph{board} y \emph{drivers} específicos, junto con los \emph{módulos sAPI}, y genera como resultado una estructura de archivos y carpetas con la biblioteca sAPI en lenguaje C para una plataforma particular. También permite generación de documentación a partir del mismo modelo.

Para su diseño se estudió una forma sistematizada de generar un módulo de código en lenguaje C. Un módulo de biblioteca en lenguaje C se compone de al menos dos archivos, uno con extensión .h (\emph{header}) que contiene declaraciones y se utiliza para incluir el código del módulo en otro archivo; y otro de extensión .c (\emph{c code}) que contiene las implementaciones. 

Para la implementación de los módulos de biblioteca sAPI se utilizan  tres archivos por módulo, un archivo \emph{module.h} que contiene la siguiente estructura:

\begin{verbatim}
  Evitar inclusión multiple{
    Inclusiones de dependencias de funciones publicas
    CPlusPlus (para cuando se usa este módulo desde C++) {
      Macros de definición de constantes publicas.
      Macros "estilo función" publicas.
      Definiciones de tipos de datos públicos.
      Prototipos de funciones publicas.
    }
  }
\end{verbatim}

Un archivo \emph{module\_private.h} con:

\begin{verbatim}
  Inclusiones de dependencias de funciones privadas.
  Macros de definición de constantes privadas.
  Macros "estilo función" privadas.
  Definiciones de tipos de datos privados.
  Prototipos de funciones privadas.
\end{verbatim}

Y un archivo \emph{module.c} con:

\begin{verbatim}
  Definiciones de variables globales públicas externas (extern).
  Definiciones de variables globales públicas.
  Definiciones de variables globales privadas (static).
  Implementaciones de funciones públicas.
  Implementaciones de funciones privadas.
\end{verbatim}

En los módulos de sAPI no se utilizan variables globales públicas y variables públicas externas para mantenerlos desacoplados.

Con esta estructura se definen las clases para formar un módulo de biblioteca en lenguaje C.

Además, se definen las clases para generar en lenguaje C todos los tipos de datos, constantes y modificadores de acceso y objetos definidos en los módulos de biblioteca sAPI.

A partir de todas estas clases se formo el generador de código en lenguaje C y documentación.

Para la implementación tanto del del generador de biblioteca sAPI como de las clases que definen el modelo se utilizó el lenguaje \emph{Javascrip} [] y el entorno de ejecución \emph{NodeJS} []. El motivo principal de esta decisión es permitir en un futuro realizar la definición de plataformas y generación de código desde una plataforma web. Estas herramientas permiten generar aplicaciones de escritorio, web o móviles con mínimos esfuerzos.

De esta forma para utilizar el generador se debe primero definir los archivos dependientes de la plataforma y luego ejecutar el generador desde una terminal con NodeJS como se esquematiza en figura \ref{fig:fullSapiGen}), que resume el producto de este trabajo final.

\begin{figure}[!htbp]
\begin{center}  % [width=14cm,height=8cm] [width=\textwidth]
\includegraphics*[width=14.5cm]{Figures/sapi_gen2.pdf}
\par\caption{Generador de la biblioteca sAPI.}\label{fig:fullSapiGen}
\end{center}
\end{figure}

Tomado los campos de descripción que incluye cada entidad del modelo genera la documentación en lenguaje \emph{Markdown}. 
 

%%----------------------------------------------------------------------------------------
%	SECTION
%----------------------------------------------------------------------------------------
\section{Análisis del software}
 
La idea de esta sección es resaltar los problemas encontrados, los criterios utilizados y la justificación de las decisiones que se hayan tomado.

Se puede agregar código o pseudocódigo dentro de un entorno lstlisting con el siguiente código:

\begin{verbatim}
\begin{lstlisting}[caption= "un epígrafe descriptivo"]

	las líneas de código irían aquí...
	
\end{lstlisting}
\end{verbatim}

A modo de ejemplo:

\begin{lstlisting}[caption=Pseudocódigo del lazo principal de control.]  % Start your code-block

#define MAX_SENSOR_NUMBER 3
#define MAX_ALARM_NUMBER  6
#define MAX_ACTUATOR_NUMBER 6

uint32_t sensorValue[MAX_SENSOR_NUMBER];		
FunctionalState alarmControl[MAX_ALARM_NUMBER];	//ENABLE or DISABLE
state_t alarmState[MAX_ALARM_NUMBER];						//ON or OFF
state_t actuatorState[MAX_ACTUATOR_NUMBER];			//ON or OFF

void vControl() {

	initGlobalVariables();
	
	period = 500 ms;
		
	while(1) {

		ticks = xTaskGetTickCount();
		
		updateSensors();
		
		updateAlarms();
		
		controlActuators();
		
		vTaskDelayUntil(&ticks, period);
	}
}
\end{lstlisting}





% Implementación
%\chapter{Implementación}
%\label{ChapterImplementacion}

%%------------------------------------------------------------------------------
%	SECTION
%------------------------------------------------------------------------------
\section{Descripción general} % Sobre la implementación.
\label{sec:implemDescripGral}

%------------------------------------------------------------------------------
%	SECTION
%------------------------------------------------------------------------------
\section{Archivos de descripción de las paltaformas del proyecto CIAA}


%%----------------------------------------------------------------------------------------
%	SECTION
%----------------------------------------------------------------------------------------
\section{Análisis del software}
 
La idea de esta sección es resaltar los problemas encontrados, los criterios utilizados y la justificación de las decisiones que se hayan tomado.

Se puede agregar código o pseudocódigo dentro de un entorno lstlisting con el siguiente código:

\begin{verbatim}
\begin{lstlisting}[caption= "un epígrafe descriptivo"]

	las líneas de código irían aquí...
	
\end{lstlisting}
\end{verbatim}

A modo de ejemplo:

\begin{lstlisting}[caption=Pseudocódigo del lazo principal de control.]  % Start your code-block

#define MAX_SENSOR_NUMBER 3
#define MAX_ALARM_NUMBER  6
#define MAX_ACTUATOR_NUMBER 6

uint32_t sensorValue[MAX_SENSOR_NUMBER];		
FunctionalState alarmControl[MAX_ALARM_NUMBER];	//ENABLE or DISABLE
state_t alarmState[MAX_ALARM_NUMBER];						//ON or OFF
state_t actuatorState[MAX_ACTUATOR_NUMBER];			//ON or OFF

void vControl() {

	initGlobalVariables();
	
	period = 500 ms;
		
	while(1) {

		ticks = xTaskGetTickCount();
		
		updateSensors();
		
		updateAlarms();
		
		controlActuators();
		
		vTaskDelayUntil(&ticks, period);
	}
}
\end{lstlisting}




%%------------------------------------------------------------------------------
\subsection{EDU-CIAA-NXP}


%%----------------------------------------------------------------------------------------
%	SECTION
%----------------------------------------------------------------------------------------
\section{Análisis del software}
 
La idea de esta sección es resaltar los problemas encontrados, los criterios utilizados y la justificación de las decisiones que se hayan tomado.

Se puede agregar código o pseudocódigo dentro de un entorno lstlisting con el siguiente código:

\begin{verbatim}
\begin{lstlisting}[caption= "un epígrafe descriptivo"]

	las líneas de código irían aquí...
	
\end{lstlisting}
\end{verbatim}

A modo de ejemplo:

\begin{lstlisting}[caption=Pseudocódigo del lazo principal de control.]  % Start your code-block

#define MAX_SENSOR_NUMBER 3
#define MAX_ALARM_NUMBER  6
#define MAX_ACTUATOR_NUMBER 6

uint32_t sensorValue[MAX_SENSOR_NUMBER];		
FunctionalState alarmControl[MAX_ALARM_NUMBER];	//ENABLE or DISABLE
state_t alarmState[MAX_ALARM_NUMBER];						//ON or OFF
state_t actuatorState[MAX_ACTUATOR_NUMBER];			//ON or OFF

void vControl() {

	initGlobalVariables();
	
	period = 500 ms;
		
	while(1) {

		ticks = xTaskGetTickCount();
		
		updateSensors();
		
		updateAlarms();
		
		controlActuators();
		
		vTaskDelayUntil(&ticks, period);
	}
}
\end{lstlisting}




%%----------------------------------------------------------------------------------------
%	SECTION
%----------------------------------------------------------------------------------------
\section{Análisis del software}
 
La idea de esta sección es resaltar los problemas encontrados, los criterios utilizados y la justificación de las decisiones que se hayan tomado.

Se puede agregar código o pseudocódigo dentro de un entorno lstlisting con el siguiente código:

\begin{verbatim}
\begin{lstlisting}[caption= "un epígrafe descriptivo"]

	las líneas de código irían aquí...
	
\end{lstlisting}
\end{verbatim}

A modo de ejemplo:

\begin{lstlisting}[caption=Pseudocódigo del lazo principal de control.]  % Start your code-block

#define MAX_SENSOR_NUMBER 3
#define MAX_ALARM_NUMBER  6
#define MAX_ACTUATOR_NUMBER 6

uint32_t sensorValue[MAX_SENSOR_NUMBER];		
FunctionalState alarmControl[MAX_ALARM_NUMBER];	//ENABLE or DISABLE
state_t alarmState[MAX_ALARM_NUMBER];						//ON or OFF
state_t actuatorState[MAX_ACTUATOR_NUMBER];			//ON or OFF

void vControl() {

	initGlobalVariables();
	
	period = 500 ms;
		
	while(1) {

		ticks = xTaskGetTickCount();
		
		updateSensors();
		
		updateAlarms();
		
		controlActuators();
		
		vTaskDelayUntil(&ticks, period);
	}
}
\end{lstlisting}





% Ensayos y Resultados
\chapter{Ensayos y Resultados}
\label{ChapterEnsayosYResultados}

\emph{Pruebas funcionales del hardware: La idea de esta sección es explicar cómo se hicieron los ensayos, qué resultados se obtuvieron y analizarlos.}


%----------------------------------------------------------------------------------------
%	SECTION
%----------------------------------------------------------------------------------------
\section{Análisis del software}
 
La idea de esta sección es resaltar los problemas encontrados, los criterios utilizados y la justificación de las decisiones que se hayan tomado.

Se puede agregar código o pseudocódigo dentro de un entorno lstlisting con el siguiente código:

\begin{verbatim}
\begin{lstlisting}[caption= "un epígrafe descriptivo"]

	las líneas de código irían aquí...
	
\end{lstlisting}
\end{verbatim}

A modo de ejemplo:

\begin{lstlisting}[caption=Pseudocódigo del lazo principal de control.]  % Start your code-block

#define MAX_SENSOR_NUMBER 3
#define MAX_ALARM_NUMBER  6
#define MAX_ACTUATOR_NUMBER 6

uint32_t sensorValue[MAX_SENSOR_NUMBER];		
FunctionalState alarmControl[MAX_ALARM_NUMBER];	//ENABLE or DISABLE
state_t alarmState[MAX_ALARM_NUMBER];						//ON or OFF
state_t actuatorState[MAX_ACTUATOR_NUMBER];			//ON or OFF

void vControl() {

	initGlobalVariables();
	
	period = 500 ms;
		
	while(1) {

		ticks = xTaskGetTickCount();
		
		updateSensors();
		
		updateAlarms();
		
		controlActuators();
		
		vTaskDelayUntil(&ticks, period);
	}
}
\end{lstlisting}





%------------------------------------------------------------------------------
%	SECTION
%------------------------------------------------------------------------------
\section{Testeo Unitario}
\label{sec:unitTest}


%------------------------------------------------------------------------------
%	SECTION
%------------------------------------------------------------------------------
\section{Banco de pruebas de hardware}
\label{sec:testBench}


%------------------------------------------------------------------------------
%	SECTION
%------------------------------------------------------------------------------
\section{Integración continua}
\label{sec:ci}


%------------------------------------------------------------------------------
%	SECTION
%------------------------------------------------------------------------------
\section{Utilización de la biblioteca para la enseñanza de programación de Sistemas Embebidos}
\label{sec:teach}


%----------------------------------------------
%
%CASOS DE USO (Paper sAPI)
%
%Desde la primer versión de la biblioteca sAPI realizada en el marco del proyecto de Java[] en 2015 se ha utilizado como ejemplo de capa de abstracción de hardware en las asignaturas dictadas por el autor. Estas asignaturas son: el  curso de posgrado “Programación de microproce- sadores”[] de la Carrera Especialización en Sistemas Embebidos (CESE[]) de la FI-UBA. Donde el autor se ha desempeñado como docente a cargo del curso durante tres ediciones, y la asignatura “Sistemas Digitales” de la carrera Ingeniería en Automatización y Control Industrial (IACI[]) de la Universidad Nacional de Quilmes (UNQ) donde el autor se desempeña en la actualidad como instructor.
%
%A partir mazo de 2016 se decide utilizar la biblioteca como en el marco de los Cursos Abiertos de Programación de Sistemas Embebidos (CAPSE[]) organizados por la ACSE[]. De esta forma se ha extendido la biblioteca de forma considerable para explicar la utilización de todos los periféricos típicos de microcontroladores con excelentes resultados en cuanto a aprendizaje por parte de los alumnos tanto de niveles avanzados como quienes dan sus primeros pasos en el aprendizaje de programación de microcontroladores.
%Además, la biblioteca sAPI se puso a disposición de cualquier persona ya que se encuentra publicada de forma libre y gratuita por internet bajo una licencia BSD modificada[] en el sitio de github del autor[] y ha sido adoptada por una gran cantidad de usuarios. 
%
%Finalmente, en diciembre de 2016 se decide utilizar la biblioteca como biblioteca estándar para las plataformas del Proyecto CIAA. Esto llevó a una profunda revisión y mejora de la misma y todavía se está trabajando en la actualidad.
%
%
%----------------------------------------------
%
%
%CASOS DE USO (paper tools)
%
%Estas herramientas se han utilizado en múltiples, entre
%ellos, el curso "Sistemas digitales" de la carrera Ingeniería en
%Automatización y Control Industrial en la UNQ, donde el
%autor actualmente se desempeña como Profesor Instructor
%desde 2014. El plan de estudio de este curso incluye
%programación avanzada en lenguaje C para
%microcontroladores, con temas tales como modularización de
%código, máquinas de estado finito, sistemas operativos en
%tiempo real que usan programación cooperativa y apropiativa
%en microcontroladores; el curso de posgrado "Programación
%de microprocesadores" dentro de la CESE de FI-UBA donde
%el autor ha sido profesor a cargo de tres ediciones. Este curso
%incluye: programación básica de microcontroladores en
%lenguaje C, modularización, máquinas de estados finitos.
%También en cursos organizados por ACSE []. En estos
%cursos las herramientas junto con la secuencia didáctica
%propuesta fueron intensamente utilizadas, mejorando esta
%secuencia entre el autor y Pablo Gómez. En los cursos de
%ACSE se distinguen dos grupos, "Cursos Abiertos de
%Programación de Sistemas Embebidos" (CAPSE), orientados a
%cualquier persona interesada en aprender la programación de
%Sistemas Embebidos; y cursos impartidos al "Instituto
%Nacional de Escuelas Técnicas" (INET) [], enfocados en la
%capacitación de docentes de Escuelas Técnicas Secundarias,
%para posteriormente poder retransmitir a sus alumnos en todo
%el país. En ambos casos, se han observado muy buenos
%resultados de aprendizaje. En el caso de los cursos CAPSE, se
%impartieron cuatro cohortes entre 2016 y 2017; en la primera
%cohorte, el 46\% de los estudiantes completaron todos los
%niveles con un promedio de 67\% de índice de aprobación por
%nivel; mientras que en la cuarta cohorte, el 73\% completó
%todos los niveles, con una tasa de aprobación promedio de
%81\% por nivel. En el caso de los cursos INET, con dos
%cohortes en 2017, el 68\% de los estudiantes completaron todos
%los niveles, con una tasa promedio de aprobación del 74\% por
%nivel. Además, después de completar sus estudios, algunos de
%estos estudiantes se inscribieron posteriormente en la CESE
%FI-UBA para continuar profundizando su aprendizaje.
%Además, se han realizado varios talleres en escuelas
%secundarias y universidades de todo el país. Vale la pena
%mencionar los múltiples tutoriales y workshop en el "Simposio
%Argentino de Sistemas Embebidos" (SASE) [].
%Todas estas herramientas son de código abierto, publicadas de
%forma gratuita en la cuenta web de github del Proyecto CIAA
%y han sido adoptadas por un número importante de usuarios y
%docentes.

%------------------------------------------------------------------------------
%	SECTION
%------------------------------------------------------------------------------
\section{Documentación y difusión}
\label{sec:documentation}


%------------------------------------------------------------------------------
\subsection{Generación automática de manual de referencia de la API en base al modelo}



%------------------------------------------------------------------------------
\subsection{Tutoriales de instalación y uso}



%------------------------------------------------------------------------------
\subsection{Difusión a la comunidad del Proyecto CIAA y Embebidos32}




% Conclusiones
\chapter{Conclusiones} % Main chapter title
\label{ChapterConclusiones} 
% Change X to a consecutive number; for referencing
% this chapter elsewhere, use \ref{ChapterX}

%----------------------------------------------------------------------------------------
%	SECTION
%----------------------------------------------------------------------------------------

\section{Trabajo realizado}

\emph{La idea de esta sección es resaltar cuáles son los principales aportes del trabajo realizado y cómo se podría continuar. Debe ser especialmente breve y concisa. Es buena idea usar un listado para enumerar los logros obtenidos.}

%----------------------------------------------------------------------------------------
%	SECTION
%----------------------------------------------------------------------------------------
\section{Próximos pasos}

\emph{Acá se indica cómo se podría continuar el trabajo más adelante.}
 

%------------------------------------------------------------------------------
%	CONTENIDO DE LA MEMORIA  - APÉNDICES
%------------------------------------------------------------------------------

% indicativo para indicarle a LaTeX los siguientes "capítulos" son apéndices
%\appendix 

% Incluir los apéndices de la memoria como archivos separadas desde la carpeta
% Appendices
% Descomentar las líneas a medida que se escriben los apéndices

%% Appendix A

\chapter{Appendix Title Here} % Main appendix title

\label{AppendixA} % For referencing this appendix elsewhere, use \ref{AppendixA}

Write your Appendix content here.
%\include{Appendices/AppendixB}
%\include{Appendices/AppendixC}

%------------------------------------------------------------------------------
%	BIBLIOGRAPHY
%------------------------------------------------------------------------------

\Urlmuskip=0mu plus 1mu\relax
\raggedright
\printbibliography[heading=bibintoc]

%------------------------------------------------------------------------------

\end{document}  
